\section{Conclusion and Future Work}
This paper presented SMAC as a set of benchmark problems for cooperative MARL.
Based on the real-time strategy game StarCraft II, SMAC focuses on decentralised
micromanagement tasks and features 14 diverse combat scenarios which challenge
MARL methods to handle partial observability and high-dimensional inputs. We
offer recommendations for reporting evaluations using standardised performance
metrics and provide a thorough report and discussion for several
state-of-the-art MARL algorithms, such as QMIX and COMA. Additionally, we are
open-sourcing PyMARL - our framework for design and analysis of deep MARL
algorithms.

In the near future, we aim to extend SMAC with new challenging scenarios that
feature a more diverse set of units and require a higher level of coordination
amongst agents. Particularly, we plan to make use of the rich skill set of
StarCraft II units, and host scenarios that require the agents to utilise the
features of the terrain. With harder multi-agent coordination problems, we aim to explore the
gaps in existing MARL approaches and motivate further research in this domain,
particularly in areas such as multi-agent exploration and coordination.

We look forward to accepting contributions from the community to SMAC and hope
that it will become a standard benchmark for measuring progress in cooperative
MARL.
